\chapter{Presets}
\label{ch:presets}

Some binding surfaces are \textbf{fixed}: they do not depend on your project at
all. Binding rosetta's own reflection API is the case that motivated this ---
the twelve \code{rosetta::script} handle classes and the eight registry free
functions are the same in every project that wants them, and copying twenty
entries into each manifest makes the three lines that \emph{are}
project-specific hard to find.

\section{Using one}

\field{preset} names a JSON fragment shipped under

\begin{center}
\code{<rosetta\_include>/rosetta/presets/<name>.json}
\end{center}

\begin{lstlisting}[style=json]
{
    "preset": "scriptable",
    "rosetta_include": "../../include",
    "user_include": ["../lib", "../lib/bindings/dynamic"],
    "user_sources": ["../lib/bindings/dynamic/auto_dynamic.cpp"],
    "module_init": { "headers": ["auto_dynamic.h"],
                     "statements": ["bank::register_all()"] },
    "targets": ["python", "lua"]
}
\end{lstlisting}

That is a complete manifest. Six keys, every one of them about \emph{this}
project: where its files are, which targets to emit, what to run at load.
\field{classes}, \field{functions} and \field{module\_name} all come from the
preset --- and a preset that supplies \field{module\_name} is what lets the
targets stay bare strings.

Give an array to apply several: \code{"preset": ["scriptable", "\ldots"]},
applied in order. An unknown name is an error naming the path that was searched,
so a typo surfaces immediately.

\section{What applying one does}

\begin{center}
\small
\begin{tabular}{@{}lL{10.0cm}@{}}
\toprule
\textbf{Merge} & Deep, with the \textbf{manifest winning} every conflict it can
  express. A key you did not write is taken from the preset; arrays present in
  both are concatenated \emph{preset first}; objects merge recursively; a scalar
  you spelled out is never overridden. \\
\addlinespace
\field{classes} / \field{functions} & Parsed with a \textbf{pristine context}:
  your \field{namespace} and \field{header\_dir} defaults describe \emph{your}
  tree and are not prepended to a preset's headers. They are appended after your
  own entries. \\
\addlinespace
\textbf{Include path} & A preset binds headers that live under
  \field{rosetta\_include}, so it puts that directory on the \textbf{target's}
  include path too. Every other binding is reflection-free by construction and
  has no reason to see rosetta's headers, which is why this is not the default. \\
\bottomrule
\end{tabular}
\end{center}

The array rule is what makes \field{module\_init} compose: the preset
contributes its caster include and \emph{your} statements survive, rather than
one silently replacing the other.

\field{classes} becomes optional when a preset supplies it --- which is why the
manifest above has none.

\section{The shipped preset}

\begin{center}
\small
\begin{tabular}{@{}lL{10.2cm}@{}}
\toprule
\textbf{Name} & \textbf{What it binds} \\
\midrule
\code{scriptable} & \code{<rosetta/script.h>} --- the bindable facade over the
  runtime object model --- plus the per-language \code{Value} casters. Lets a
  host language walk the metadata of any library that emits \lang{dynamic}
  tables, and call into it by name. \\
\bottomrule
\end{tabular}
\end{center}

Chapter~\ref{ch:scriptable} is what that preset is \emph{for}, and is worth
reading before using it --- it is the most unusual thing in rosetta and the
hardest to guess at from the manifest alone.

\section{Writing your own}

A preset is only a JSON fragment. Drop one into
\code{<rosetta\_include>/rosetta/presets/} and it is available by name to every
manifest that points at that include directory.

It is worth doing when a binding surface is genuinely fixed across projects:

\begin{itemize}
  \item a company-wide utility library every binding also exposes;
  \item a facade type plus its helpers, where the twenty entries never vary;
  \item a set of \field{sequences} / \field{matrices} / \field{interop}
    registrations shared by every project using the same numerical stack.
\end{itemize}

It is \emph{not} worth doing to factor out the parts of a manifest that differ
per project --- \field{user\_include}, \field{user\_sources},
\field{targets} --- because the merge rule (manifest wins) means those would be
overridden anyway, and a preset nobody can use unchanged is just an indirection.

The fields most worth putting in one:

\begin{lstlisting}[style=json]
{
  "module_name": "mylib",
  "classes":   [ ... fixed entries ... ],
  "functions": [ ... fixed entries ... ],
  "module_init": { "headers": ["mylib/init.h"],
                   "statements": ["mylib::initialize()"] },
  "sequences": ["mylib::Vector"]
}
\end{lstlisting}

\begin{note}
Remember the pristine-context rule when writing one: a preset's
\field{classes} entries must spell their headers and names in full, because
the consuming manifest's \field{namespace} and \field{header\_dir} will
\emph{not} be applied to them.
\end{note}
